% from ega4-18-fr.tex
\oldpage[IV-4]{220}

\emph{Dans ces conditions:}
\begin{enumerate}
  \item[(i)] Supposons $S_\alpha$ quasi-compact. Soit $\sh F_\alpha$ un $\sh O_{S_\alpha}$-Module quasi-cohérent, que l'on suppose en outre de type fini quand l'hypothèse $1^\circ$ est vérifiée. Soit $(f_{i\alpha})_{1\leq i\leq n}$ une suite finie de sections de $\sh F_\alpha$ au-dessus de $S_\alpha$. Pour que la suite des sections $f_i$ du $\sh O_S$-Module $\sh F$ au-dessus de $S$, correspondant aux $f_{i\alpha}$ $(1\leq i\leq n)$, soit $\sh F$-régulière, il faut et il suffit qu'il existe $\lambda\geq\alpha$ tel que la suite des sections $f_{i\lambda}$ de $\sh F_\lambda$ au-dessus de $S_\lambda$ soit $\sh F_\lambda$-régulière.
  \item[(ii)] Supposons $X_\alpha$ quasi-compact, et soit $j_\alpha:X_\alpha\to S_\alpha$ une immersion, que l'on suppose localement de présentation finie lorsque l'hypothèse $2^\circ$ est vérifiée. Alors, pour que l'immersion correspondante $j:X\to S$ soit régulière, il faut et il suffit qu'il existe $\lambda\geq\alpha$ tel que $j_\lambda:X_\lambda\to S_\lambda$ soit régulière.
\end{enumerate}
\end{proposition}

\begin{enumerate}
  \item[(i)] Si $p_\lambda:S\to S_\lambda$ est la projection canonique, on a
  \[
    \sh F/\Bigl(\sum_{j=1}^{i-1}f_j\sh F\Bigr)
    =p_\lambda^*\Bigl(\sh F_\lambda/\Bigl(\sum_{j=1}^{i-1}f_{j\lambda}\sh F_\lambda\Bigr)\Bigr),
  \]
  donc on est ramené au cas $n=1$, auquel cas on supprime l'indice $i$. Le fait que la condition soit suffisante résulte de ce que $p_\lambda$ est plat (8.3.8) et de ($0$, 15.1.5). Dans le cas $2^\circ$, $p_\lambda$ est fidèlement plat (8.3.8) et la nécessité de la condition résulte encore de ($0$, 15.2.5), avec $\lambda=\alpha$. Dans le cas $1^\circ$, notons $\sh N_\lambda$ (resp. $\sh N$) le noyau de l'homomorphisme $\sh F_\lambda\to\sh F_\lambda$ (resp. $\sh F\to\sh F$), multiplication par $f_\lambda$ (resp. par $f$); puisque $\sh F_\lambda$ est cohérent par hypothèse, il en est de même de $\sh N_\lambda$, donc l'hypothèse $\sh N=0$ entraîne $\sh N_\lambda=0$ pour un $\lambda\geq\alpha$ en vertu de (8.5.8, (ii)).
  \item[(ii)] Comme $p_\lambda$ est plat, la suffisance de la condition résulte de (19.1.5, (ii)). Pour prouver sa nécessité, notons que, puisque $j_\alpha(X_\alpha)$ est quasi-compact, il est contenu dans un ouvert quasi-compact de $S_\alpha$ et on peut donc se limiter au cas où $S_\alpha$ est aussi quasi-compact et $j_\alpha$ une immersion fermée, de sorte que l'image de $X_\alpha$ est définie par un Idéal quasi-cohérent $\sh J_\alpha$ de $\sh O_{S_\alpha}$, qu'on peut en outre supposer de type fini dans les cas $1^\circ$ et $2^\circ$ ($j_\alpha$ étant localement de présentation finie dans le cas $2^\circ$); l'image de $X_\lambda$ (resp. $X$) dans $S_\lambda$ (resp. $S$) est alors définie par $\sh J_\lambda=\sh J_\alpha\sh O_{S_\lambda}$ (resp. $\sh J=\sh J_\alpha\sh O_S$) qui est encore de type fini. On peut en outre, compte tenu de (8.2.11), supposer que $\sh J$ est engendré par une suite régulière $(f_i)_{1\leq i\leq n}$ de sections au-dessus de $S$, qui définissent par suite un homomorphisme surjectif $u:\sh O_S^n\to\sh J$. Compte tenu de (8.5.2, (i)) et (8.5.7), il existe un $\lambda\geq\alpha$ et un homomorphisme surjectif $u_\lambda:\sh O_{S_\lambda}^n\to\sh J_\lambda$ tels que $u=p_\lambda^*(u_\lambda)$, donc les $f_i$ sont les images canoniques de sections $f_{i\lambda}$ de $\sh J_\lambda$ au-dessus de $S_\lambda$ engendrant cet Idéal. En vertu de (i), il existe donc $\mu\geq\lambda$ tel que la suite $(f_{i\mu})$ soit $\sh O_{S_\mu}$-régulière, donc l'immersion $j_\mu$ est régulière.
\end{enumerate}

\begin{proposition}[19.8.2]
\label{IV.19.8.2}
Soient $X_\alpha$ un $S_\alpha$-préschéma quasi-compact et localement de présentation finie.
\begin{enumerate}
  \item[(i)] Soient $\sh F_\alpha$ un $\sh O_{X_\alpha}$-Module quasi-cohérent de présentation finie, $(f_{i\alpha})_{1\leq i\leq n}$ une suite de sections de $\sh O_{X_\alpha}$ au-dessus de $X_\alpha$. Pour que la suite de sections $f_i$ de $\sh O_X$ au-dessus de $X$, correspondant aux $f_{i\alpha}$, soit $\sh F$-transversalement régulière relativement à $S$ (19.2.1), il faut et il suffit qu'il existe $\lambda\geq\alpha$ tel que la suite des sections $f_{i\lambda}$ de $\sh O_{X_\lambda}$ au-dessus de $X_\lambda$ soit $\sh F_\lambda$-transversalement régulière relativement à $S_\lambda$.
  \item[(ii)] Soient $Y_\alpha$ un $S_\alpha$-préschéma quasi-compact, $j_\alpha:Y_\alpha\to X_\alpha$ une $S_\alpha$-immersion localement

